\documentclass{standalone}

\usepackage{tikz}
\usepackage{pgfplots}
\usepackage{amsmath}

\usetikzlibrary{arrows.meta,fadings,decorations.pathmorphing}

\makeatletter
\newif\iftikz@shading@path
\tikzset{
    % Need this for the bounding box calculation of the
    % fading and the rectangle that bounds the fading.
    % If set to zero things (e.g., arrows) may get cut off.
    fading sep/.store in=\fadingsep,
    fading sep=0.25cm,
    shaded path/.code={%
        % Prevent this stuff happning recursively.
        \iftikz@shading@path%
        \else%
            \tikz@shading@pathtrue%
            \tikz@addmode{%
                % Interrupt the picture to create a fading.
                \pgfinterruptpicture%
                    \begin{tikzfadingfrompicture}[name=.]
                        \path [shade=none,fill=none]#1;%
                        % Need to set the bounding box manually. Include the \fadingsep border.
                        \xdef\fadingboundingbox{{\noexpand\pgfpoint{\the\pgf@picminx-\fadingsep}{\the\pgf@picminy-\fadingsep}}%
                            {\noexpand\pgfpoint{\the\pgf@picmaxx+\fadingsep}{\the\pgf@picmaxy+\fadingsep}}}%
                        \expandafter\pgfpathrectanglecorners\fadingboundingbox%
                        \pgfusepath{discard}%
                    \end{tikzfadingfrompicture}%
                \endpgfinterruptpicture%
                % Install a rectangle that covers the shaded/faded path.
                \expandafter\pgfpathrectanglecorners\fadingboundingbox%
                % Make the fading happen.
                \def\tikz@path@fading{.}%
                \tikz@mode@fade@pathtrue%
                \tikz@fade@adjustfalse%
                \pgfpointscale{0.5}{\expandafter\pgfpointadd\fadingboundingbox}%
                \def\tikz@fade@transform{shift={(\the\pgf@x,\the\pgf@y)}}%
            }%
        \fi%
    }
}

\begin{document}
\begin{tikzpicture}
    \begin{axis}[ylabel=\Large{$T/T_c$}, xlabel=\Large{$z/\ell$}, legend pos=north west]
        \foreach[evaluate=\n as \redfrac using (\n-1)*100/(20-1)] \n in {2,3,...,20}{
            \edef\temp{\noexpand\addplot[color=blue!\redfrac!red] table[col sep=comma, x index=0, y index=\n]{temperature_profile.csv};}
            \temp
        }
        \addplot[domain=0:1, samples=5, color=red, no marks] coordinates{(0,1) (0,10) (1,10)};

        \path [left color=red, right color=blue, shaded path={ 
    [draw=transparent!0, thick, ->] (axis cs:0.1,9)
        to[out=-20,in=110]
    (axis cs:0.65,1.7);
}];

        % \draw[very thick,->, ]
        % (axis description cs:0.05,0.95)
        %     to[out=-20,in=130]
        % (axis description cs:0.95,0.05);
    \end{axis}
\end{tikzpicture}
\end{document}